\chapter{Những câu hỏi làm thầy khựng lại}
\markboth{Những câu hỏi làm thầy khựng lại}{Những câu hỏi làm thầy khựng lại}

\section{Người lớn sai nhiều vậy sao cứ bắt tụi em chuẩn}
\begin{truyen}{Người lớn sai nhiều vậy sao cứ bắt tụi em chuẩn}{Classroom Talk}
\chuhoa{M}{ột bạn hỏi câu làm lớp im: ``Người lớn cũng nóng tính, hứa rồi quên, nói một đằng làm một nẻo. Vậy sao cứ bắt tụi em phải chuẩn?''}

Thầy đứng yên vài giây rồi đáp: ``Câu đó không sai. Người lớn nhiều khi chưa làm gương đủ tốt. Nhưng cái sai của người lớn không làm cái đúng bớt đúng.''

Bạn ấy hỏi tiếp: ``Vậy tụi em phải nhìn vô ai?''

Thầy nói: ``Nhìn vào điều đúng, rồi cũng nhìn kỹ để sau này mình đỡ lặp lại cái dở của người lớn.''

\textit{(Adult inconsistency is real, but it does not erase the value of what is still right.)}
\end{truyen}

\begin{baihoc}
Thấy người lớn sai không phải để khỏi phải tốt hơn. Là để biết mình cần sống tử tế hơn ở chỗ nào.
\textit{(Seeing adult flaws should not excuse your own. It should sharpen your sense of what to improve.)}
\end{baihoc}

\begin{ghinhoanh}
\textbf{Short English to remember:}

Adult flaws are real.

Truth does not disappear because messengers fail.
\end{ghinhoanh}

\begin{tuhocnhanh}
1. Em có đang lấy cái sai của người khác để nhẹ lỗi cho mình không? \textit{(Do you use someone else’s wrong as a discount for your own?)}

2. Điều tốt nào em vẫn muốn giữ dù từng thấy người lớn làm chưa tới? \textit{(What good value do you still want to keep even after seeing adults fail at it?)}
\end{tuhocnhanh}
\ngancach

\section{Thầy từng bất lực muốn mặc kệ chưa}
\begin{truyen}{Thầy từng bất lực muốn mặc kệ chưa}{Classroom Talk}
\chuhoa{C}{ó bạn hỏi nhỏ nhưng thẳng: ``Thầy có bao giờ dạy một lớp xong mệt quá, thấy tụi em lì quá rồi chỉ muốn mặc kệ không?''}

Thầy cười buồn: ``Có chứ. Giáo viên cũng là người. Có lúc nói mấy lần mà không vào, lòng nản là thật.''

Bạn ấy hỏi tiếp: ``Vậy sao thầy còn làm?''

Thầy đáp: ``Vì nếu thầy mặc kệ luôn, nhiều đứa sẽ quen với cảm giác không còn ai kỳ vọng ở mình nữa.''

\textit{(Teachers also feel exhaustion and helplessness, but staying engaged can matter deeply to students who are close to giving up.)}
\end{truyen}

\begin{baihoc}
Sự kiên nhẫn của người dạy không vô hạn. Nhưng đôi khi nó cứu được một đứa đang sắp buông mình.
\textit{(A teacher’s patience is not unlimited, but it can sometimes rescue a student on the edge of quitting.)}
\end{baihoc}

\begin{ghinhoanh}
\textbf{Short English to remember:}

Teachers get tired too.

Staying can still change a life.
\end{ghinhoanh}

\begin{tuhocnhanh}
1. Em có từng làm ai muốn bỏ cuộc với mình chưa? \textit{(Have you ever made someone feel like giving up on you?)}

2. Nếu có người vẫn chưa bỏ em, em nên đáp lại thế nào? \textit{(If someone has not given up on you, how should you respond?)}
\end{tuhocnhanh}
\ngancach

\section{Học sinh cần bị quản hay cần được hiểu}
\begin{truyen}{Học sinh cần bị quản hay cần được hiểu}{Classroom Talk}
\chuhoa{M}{ột bạn hỏi: ``Nhiều khi tụi em bị quản kỹ quá thấy ngộp. Nhưng nếu thả lỏng lại dễ loạn. Vậy tụi em cần bị quản hay cần được hiểu hơn?''}

Thầy đáp: ``Cả hai, nhưng theo thứ tự đúng. Hiểu để biết phải quản chỗ nào. Quản để giữ cho việc hiểu không thành nuông chiều.''

Bạn ấy gật gù: ``Nghĩa là thương không đồng nghĩa với dễ dãi?''

Thầy nói: ``Đúng. Có khi nghiêm mới là một dạng thương có trách nhiệm.''

\textit{(Students need understanding, but understanding without structure can become harmful indulgence.)}
\end{truyen}

\begin{baihoc}
Thầy cô tốt không phải lúc nào cũng mềm. Nhiều khi họ đang giữ cho mình khỏi trượt dài.
\textit{(Good teachers are not always gentle in appearance. Sometimes they are preventing a longer fall.)}
\end{baihoc}

\begin{ghinhoanh}
\textbf{Short English to remember:}

Understanding matters.

Structure matters too.
\end{ghinhoanh}

\begin{tuhocnhanh}
1. Em đang cần được cảm thông hay đang cần bị kéo lại cho đúng? \textit{(Do you need compassion, or do you need correction right now?)}

2. Em có nhầm giữa được hiểu và được chiều không? \textit{(Do you confuse being understood with being indulged?)}
\end{tuhocnhanh}
\ngancach

\section{Có đứa im không phải ngoan, mà là hết muốn nói}
\begin{truyen}{Có đứa im không phải ngoan, mà là hết muốn nói}{Classroom Talk}
\chuhoa{C}{ó bạn nhìn về cuối lớp rồi hỏi: ``Mấy bạn ít nói, ít phá có chắc là ổn không thầy? Hay có đứa im vì chẳng còn muốn mở miệng nữa?''}

Thầy chậm giọng: ``Em hỏi đúng chỗ đau của nhiều lớp. Không phải học sinh im nào cũng ngoan. Có em im vì sợ, vì mệt, vì thấy nói cũng chẳng ai nghe.''

Bạn ấy hỏi tiếp: ``Vậy thầy cô làm sao biết?''

Thầy đáp: ``Phải nhìn lâu hơn điểm số và hành vi bề mặt. Nhiều đứa ngoan quá trên giấy nhưng đang mòn dần ở bên trong.''

\textit{(Quiet students are not always fine. Silence can hide exhaustion, fear, or withdrawal.)}
\end{truyen}

\begin{baihoc}
Đừng chỉ chú ý những đứa ồn. Nhiều khi đứa cần cứu nhất lại là đứa im nhất.
\textit{(Do not only notice the loud students. The quietest one may need help the most.)}
\end{baihoc}

\begin{ghinhoanh}
\textbf{Short English to remember:}

Silence is not always peace.

Quiet students need attention too.
\end{ghinhoanh}

\begin{tuhocnhanh}
1. Em có đang im vì bình an hay vì cạn năng lượng? \textit{(Are you quiet because you are calm, or because you are drained?)}

2. Em có người lớn nào em có thể nói thật không? \textit{(Do you have an adult you can speak honestly with?)}
\end{tuhocnhanh}
\ngancach